Assessing a fitted field also requires distinguishing parameter response from observational agreement and physical feasibility. Analytical sensitivity calculations already quantify how laser inputs and constant thermophysical properties change pool dimensions and temperature histories~\cite{yang2024}. These responses describe the adopted model; comparison with observations determines which outputs it can support. Combined thermography and metallography of 316L, IN625 and CoCr tracks found useful surface-geometry trends, while strong depth correlations coexisted with systematic deviations outside stable conduction melting~\cite{soares2026}. This separates correlation from dimensional agreement and makes the operating regime part of the assessment. A complementary depth-based test inverts a specified conduction model for the absorptivity needed to reproduce measured penetration, then compares that value with its physical bound and tests alternative effective-transport assumptions~\cite{hong2026}. The recovered factor serves as an attainability diagnostic under specified assumptions. Together, these studies distinguish sensitivity, output agreement and attainability, providing a basis for examining the constitutive assumptions that shape each modeled observable.
